% Reorganized shared appendix file; load once in the documented order.
\section{Positive systems, siphons, triangularity of the Jacobian on siphon faces and the Metzler property of its transversal block, and invasion exponents}\label{ss:siphons}

\beD[Positive ODE][][pos]
An ODE $\dot x=f(x)$, with $f$ locally Lipschitz on a neighbourhood of
$\R^{n_S}_{\ge0}$, is called \emph{positive} \cite{rantzer2015scalable} or
\emph{non-negative} \cite{haddad2010} if the orthant $\R^{n_S}_{\ge0}$ is forward
invariant. Equivalently, $f_i(x)\ge0$ whenever $x\ge0$ and $x_i=0$.
\eeD

\beD[Siphon and minimal siphon][siphon][Siphon]
For a set $\Sigma$ of coordinates, let
$\mathcal F_\Sigma=\{x\in\R^{n_S}_{\ge0}:\ x_i=0\ \text{for all } i\in\Sigma\}$
be the corresponding face of the orthant. A nonempty set $\Sigma$ is a
\emph{siphon} of a positive ODE if $\mathcal F_\Sigma$ is forward invariant, and
a \emph{minimal siphon} if no nonempty proper subset of $\Sigma$ is a siphon.
\eeD

Since the coordinates in $\Sigma$ vanish on $\mathcal F_\Sigma$, the face is
forward invariant if and only if $f_i$ vanishes on $\mathcal F_\Sigma$ for
every $i\in\Sigma$. A siphon is thus a set of variables that cannot be created
from nothing: once absent, they stay absent. For a reaction network with
mass-action kinetics and positive rate constants, or more generally with nonnegative rates that vanish when a reactant is absent
and are positive when all reactants are present, this is the
combinatorial condition that every reaction producing a species of $\Sigma$
consumes one; for more general rates, additional invariant faces can arise from
rates that vanish for other reasons.

\beD[Exact zero set, or exact supra-siphon, of a point][exact supra-siphon][ExZ]
The \emph{exact zero set} of a point $w\in\R^{n_S}_{\ge0}$ is $Z(w)=\{i:\ w_i=0\}$. When $Z(w)$ is
a siphon it is also called the \emph{exact supra-siphon} of $w$. Since $w\in\mathcal F_\Sigma$
exactly when $\Sigma\subseteq Z(w)$, the exact supra-siphon is the largest siphon on whose face
$w$ lies, it contains every siphon whose face contains $w$, and $w$ lies in the relative interior of that
siphon face,
\[
\mathcal F^\circ_{Z(w)}=\{x\in\mathcal F_{Z(w)}:\ x_j>0\ \text{for } j\notin Z(w)\}.
\]
For $w\gg0$, $Z(w)=\varnothing$ and $w$ lies in the interior of the orthant.
\eeD



\beP[The exact zero set of a RN boundary omega-limit point is a siphon][exact zero sets are siphons][zeroset]\label{p:AdLS}
Consider the reaction system
\[
\dot x=f(x)=\sum_r v_r(x)(b_r-a_r),
\]
where $a_r,b_r\in\R_{\ge0}^{n_S}$ are the reactant and product vectors. Assume that the rates are $C^1$
on a neighbourhood of the nonnegative orthant and, for $x\ge0$, satisfy
\[
v_r(x)\ge0,\qquad
v_r(x)>0
\quad\Longleftrightarrow\quad
x_j>0\ \text{for every }j\text{ with }a_{rj}>0
\]
(these assumptions hold for mass-action kinetics with positive rate constants). Let $x(t)$ be a
forward-complete solution with $x(0)\gg0$, and let $w$ be a boundary $\omega$-limit point. Then
its exact zero set $Z(w)$ is a siphon: every reaction producing a species in $Z(w)$ consumes at
least one species in $Z(w)$. The same holds for every boundary equilibrium $w$.
\eeP
\begin{proof}
First we show that $f_i(w)=0$ for every $i\in Z(w)$; for an equilibrium this holds by assumption.
Choose $t_k\to\infty$ with $x(t_k)\to w$. Local existence, uniqueness and continuous dependence
give a solution $y:(-\delta,\delta)\to\R^{n_S}$ through $y(0)=w$ such that, for each
$s\in(-\delta,\delta)$,
\[
y(s)=\lim_{k\to\infty}x(t_k+s)\ge0.
\]
Thus, for $i\in Z(w)$, the differentiable function $y_i$ has a local minimum at $s=0$, and
$f_i(w)=y_i'(0)=0$.
Fix $i\in Z(w)$. Every reaction consuming species $i$ has zero rate at $w$, because $w_i=0$. Hence
\[
0=f_i(w)=\sum_{r:\,a_{ri}=0}b_{ri}v_r(w).
\]
All summands are nonnegative. Therefore every reaction producing $i$ without consuming it has zero
rate at $w$. By the rate assumption, such a reaction consumes some species whose concentration at $w$
is zero, hence a species in $Z(w)$. A reaction that consumes $i$ already consumes a species in
$Z(w)$. This proves the siphon property.
\end{proof}

\beXa[For a general positive ODE, an exact zero-set of  a boundary omega-limit point  may  fail to be siphon][][ZsNsip] On $[0,1]\times[0,2]$ the system
$\dot x=(y-1)^2-x$, $\dot y=1-y$ sends every trajectory to $(0,1)$, whose exact zero set $\{x\}$ is not a siphon,
since $\dot x=(y-1)^2>0$ on $\{x=0\}$ for $y\ne1$. The exact zero set of a point is distinguished from the larger invariant
faces containing it.
\eeXa

For a   siphon $\Sigma$, $z\in\R^m_{\ge0}$ collects the coordinates in
$\Sigma$ and $y\in\R^n_{\ge0}$ the remaining, \emph{resident}, ones; we write
$\dot z=f_z(z,y)$, $\dot y=f_y(z,y)$, with $f$ of class $C^1$, and at a point
$(0,y)$ of the face we call
\[
\Jtr(y):=D_zf_z(0,y),\qquad J_{\mathrm{tan}}(y):=D_yf_y(0,y),\qquad
J_{\mathrm{mix}}(y):=D_zf_y(0,y)
\]
the \emph{transversal}, \emph{tangential} and \emph{mixed} blocks of the
Jacobian. % a driver that redefines \Jtr as J_tra gets the short wording
\providecommand{\SiInJdefault}{J^\perp}\ifx\Jtr\SiInJdefault The transversal block is also written $J^\perp$ or
$\JtrS{\Sigma}$; these are the same matrix.\else When the siphon is to be indicated, the transversal block is
written $\JtrS{\Sigma}$.\fi{} Invariance of the face reads
\begin{equation}\label{eq:faceinv}
f_z(0,y)=0\qquad\text{for all } y\in\R^n_{\ge0}.
\end{equation} $J_{\mathrm{tan}}$ is the Jacobian of the ODE restricted to the face,
while $\Jtr$ governs the linearized dynamics of the absent variables.



Below, $s(A)=\max\{\operatorname{Re}\mu:\ \mu\in\operatorname{spec}A\}$ denotes the
spectral abscissa of a square matrix $A$.

\beP[The Jacobian on a siphon face is block triangular, with a Metzler transversal block][][Tri]
\label{p:Sta}
Let $\Sigma$ be a siphon of a positive $C^1$ system. Then:
\BEN
\im At every point of its face, $D_yf_z(0,y)=0$ and
\begin{equation}\label{eq:blocks}
J(0,y)=
\begin{pmatrix}
\Jtr(y)&0\\
J_{\mathrm{mix}}(y)&J_{\mathrm{tan}}(y)
\end{pmatrix}.
\end{equation}
Its spectrum is the union of the spectra of the diagonal blocks.
\im An equilibrium $E=(0,y^*)$ is locally exponentially stable if and only if
$J_{\mathrm{tan}}(y^*)$ and $\Jtr(y^*)$ are Hurwitz. It is unstable if either
block has an eigenvalue with positive real part.
\im The transversal block $\Jtr(y)$ is Metzler at every point of the face;
the tangential block need not be Metzler.
\EEN
The same conclusions hold on a face of an invariant box $[0,\bar x]$, $\bar x\gg0$,
with derivatives understood in admissible directions when the vector field
is only $C^1$ up to the boundary. Empty blocks impose no condition.
\eeP
\begin{proof}
Differentiating~\eqref{eq:faceinv} in the resident variables gives
$D_yf_z(0,y)=0$ in the relative interior and, by continuity, on the entire
face. The characteristic polynomial therefore factors over the two diagonal
blocks, giving the stability and instability assertions. Finally, for distinct $i,j\in\Sigma$, positivity at $(te_j,y)$ gives
\[
(f_z)_i(te_j,y)\ge0=(f_z)_i(0,y)\qquad(t>0).
\]
Dividing by $t$ and letting $t\downarrow0$ proves
$(\Jtr)_{ij}(y)\ge0$. These arguments use only admissible directions and
apply equally on a box.
\end{proof}


\beD[Transversal block and invasion exponent][invasion exponent][Inv]
At a fixed point $E=(0,y^*)$ of the face of a siphon $\Sigma$, the
\emph{transversal block} is $M_\Sigma(E):=\Jtr(y^*)$, and the \emph{invasion
exponent} of $\Sigma$ at $E$ is
\[
\lambda_\Sigma(E):=s\bigl(M_\Sigma(E)\bigr).
\]
\eeD

By Proposition~\ref{p:Tri}, $M_\Sigma(E)$ is Metzler and
$\lambda_\Sigma(E)$ is an eigenvalue of the full Jacobian. For sufficiently
large $c$, Perron--Frobenius theory gives
$\lambda_\Sigma(E)=\rho(M_\Sigma(E)+cI)-c$, with a nonnegative eigenvector.
It measures asymptotic growth in the transverse linearized system:
$\lambda_\Sigma(E)>0$ permits invasion, whereas
$\lambda_\Sigma(E)<0$ gives exponential decay of all transverse perturbations;
at equality, linearization does not decide invasion.

\beP[transverse factorization on siphon faces][][traFac]\label{t:Inv}
For a reaction network  as defined in Proposition \ref{p:zeroset}, write $\nu^-_\ell\to\nu^+_\ell$, and   assume
each rate vanishes exactly when at least one of its reactant
concentrations vanishes (this  holds for mass-action
kinetics).

For any siphon $\Sigma$
\[
 f_\Sigma(z,y)=M_\Sigma(z,y)z,\qquad
 M_\Sigma(z,y)=\int_0^1D_zf_\Sigma(tz,y)\,dt,
 \qquad M_\Sigma(0,y)=\JtrS{\Sigma}(y).
\]
The matrix $M_\Sigma$ need not be Metzler away from the face. If the ODE
preserves the nonnegative orthant, the block-triangular Jacobian on
$\mathcal F_\Sigma$, its Metzler transverse block, and the stability test at a face
equilibrium are those of Proposition~\ref{p:Tri}.
\eeP
\begin{proof}
On the face all consumption rates of absent species vanish. If the
reaction condition holds, all production rates vanish too. Conversely,
at a relative-interior point of the face, a reaction producing an
absent species without an absent reactant has positive rate. It cannot
be cancelled by consumption of that species, so the face is not
invariant. The integral identity is the fundamental theorem of
calculus applied along $t\mapsto(tz,y)$; at $z=0$ it gives
$M_\Sigma(0,y)=D_zf_\Sigma(0,y)=\JtrS{\Sigma}(y)$.
\end{proof}

\paragraph{Siphons in epidemiology.}
In the multi-strain models considered here, \emph{epidemiologic strains} often
correspond to \emph{minimal siphons}. By a \emph{mutation} we mean only a coupling
between strains that changes the Perron--Frobenius form of the transversal block
from block diagonal to block triangular. The totality of ``infected variables'' is
the total siphon, the union of all minimal siphons, and a \DFE\ (DFE) lies on its
face; a larger siphon may be needed when the face of the total siphon does not
contain a unique fixed point. For a splitting $\Jtr=F-V$ with
$F\ge0$ and $V$ a nonsingular $M$-matrix \cite{Diek,Van},
the Perron spectral radius $\rho(FV^{-1})$ of the next generation matrix on a
siphon, expressed as a function of the resident variables and called a
reproduction function in \cite{AAN,AAHJ,AABH25}, characterizes invasion of boundary
fixed points\bookref{p:NGM}{ (Proposition~\ref{p:NGM})}{}, and often their existence domains; other eigenvalues of the next
generation matrix play no role in the linear invasion threshold. \bookref{p:toolExit}{The admissibility of a biologically specified $V$ is tested in
Theorem~\ref{p:toolExit}.}{}

% Exact-support counting: unused by the short paper, kept for the long versions.
\ifshort\else
\subsection{Equilibria are partitioned by their exact supra-siphons}
Counting equilibria, and counting the stable ones, needs two ingredients from
above: the exact zero set (Definition~\ref{d:ExZ} and
Proposition~\ref{p:zeroset}), which says \emph{where} an equilibrium is counted,
and the block-triangular Jacobian (Proposition~\ref{p:Tri}), which says
\emph{whether} it is stable. Throughout, $\dot x=f(x;\theta)$ depends on
parameters $\theta$ ranging over a region on which the family of siphons does
not change.

\beD[Exact-support counts and the full-stability weight][full-stability weight][FSW]
For a siphon $\Sigma$, or for $\Sigma=\varnothing$, write
\[
\mathcal F_\Sigma^\circ=\{x\in \mathcal F_\Sigma:\ x_j>0\ \text{for } j\notin \Sigma\},
\]
the interior of the siphon face $\mathcal F_\Sigma$; $\mathcal F_\varnothing^\circ$ is the interior
of the orthant, and $\varnothing$ is used only as an index, not as a siphon.
The points of $\mathcal F_\Sigma^\circ$ are exactly the points whose exact zero set is $\Sigma$.
Let
\begin{enumerate}
\item $N_\Sigma(\theta)$ be the number of equilibria in $\mathcal F_\Sigma^\circ$, and
$N_\Sigma^{\mathrm{st}}(\theta)$ the number of those that are locally exponentially
stable;
\item $N_{\ge0}(\theta)$ and $N_{\ge0}^{\mathrm{st}}(\theta)$ be the number of
all nonnegative equilibria and of the locally exponentially stable ones.
\end{enumerate}
At a point $(0,y)\in \mathcal F_\Sigma^\circ$, with the blocks of Proposition~\ref{p:Tri}
for the siphon $\Sigma$, the \emph{full-stability weight} is
\begin{equation}\label{eq:weight}
W_\Sigma(y;\theta)=
\ind{J_{\mathrm{tan},\Sigma}(y;\theta)\text{ Hurwitz}}\,
\ind{\JtrS{\Sigma}(y;\theta)\text{ Hurwitz}},
\end{equation}
where an empty block counts as Hurwitz. The first factor tests stability
within the siphon face, the second tests that the absent variables cannot
invade.
\eeD

\beP[The exact supra-siphons partition the equilibria, and the full-stability weight selects the stable ones][counting by exact supra-siphon][Part]
Assume that every boundary equilibrium has a siphon as its exact zero set; by
Proposition~\ref{p:zeroset} this holds for reaction networks under its rate
assumptions. Then:
\begin{enumerate}
\item every nonnegative equilibrium lies in exactly one set $\mathcal F_\Sigma^\circ$, namely
the one indexed by its exact supra-siphon, or by $\varnothing$ if it is
positive; consequently
\[
N_{\ge0}=\sum_\Sigma N_\Sigma,\qquad N_{\ge0}^{\mathrm{st}}=\sum_\Sigma N_\Sigma^{\mathrm{st}},
\]
the sums running over $\varnothing$ and all siphons, not only the minimal
ones;
\item an equilibrium $(0,y)\in \mathcal F_\Sigma^\circ$ is locally exponentially stable if
and only if $W_\Sigma(y;\theta)=1$, so that
$N_\Sigma^{\mathrm{st}}=\sum_{(0,y)\in \mathcal F_\Sigma^\circ,\ f(0,y)=0}W_\Sigma(y;\theta)$.
\end{enumerate}
\eeP
\begin{proof}
For item~1, an equilibrium $e$ lies in $\mathcal F_\Sigma^\circ$ exactly when $\Sigma=Z(e)$, and
the sets $\mathcal F_\Sigma^\circ$ for distinct $\Sigma$ are therefore disjoint; by assumption
$Z(e)$ is a siphon whenever $e$ is on the boundary. Item~2 is item~2 of
Proposition~\ref{p:Tri} applied to the siphon $\Sigma=Z(e)$.
\end{proof}

\beR[Why minimal siphons cannot be used as counting classes][][noDouble]
Summing over the minimal siphons instead of the exact supra-siphons counts some
equilibria several times when minimal siphons overlap. For example, with the
two overlapping minimal siphons $\{1,2\}$ and $\{2,3\}$, the origin lies on the
faces of both, but its exact supra-siphon is $\{1,2,3\}$, so it is counted once.
Likewise, the first factor of $W_\Sigma$ alone counts the attractors of the dynamics
restricted to the siphon face, some of which are invaded by the absent
variables. Without the assumption of Proposition~\ref{p:Part}, a boundary
equilibrium may lie in no counting class at all: in
Example~\ref{e:ZsNsip} the exact zero set $\{x\}$ of the limit point is not a
siphon.
\eeR
\bookref{t:Ric}{These counts are used in Theorem~\ref{t:Ric}, whose stability
weight is \eqref{eq:weight}, and in the sum~\eqref{eq:sum} over exact
supra-siphons.}{}
\fi



\beR[historical notes][][hisS]
The invariant face-siphons correspondence and its use in persistence are developed in
\cite{AdLS,ShiuStu}.

The Jacobian and transversal factorizations are  elementary \cite{AAHK}, \cite{AAH26} and also supply
the idea of the role of Perron data in face-repulsion.

The exact-zero-set assertion of Proposition~\ref{p:zeroset} is already contained
in \cite[Lemma~5.3 and Proposition~5.4]{AdLS}, for real-analytic rates; the proof above gives a
short argument using a local trajectory through the omega-limit point, and needs only $C^1$ rates.
The invariant-face equivalence is \cite[Proposition~5.5]{AdLS}.
\eeR 